\subsection{EXTENSION OF UNIFORMLY CONTINUOUS FUNCTIONS}

The theorem of extension by continuity (Chapter I, § 8, no. 5, Theorem 1) has important additions when the functions in question take their values in a complete Hausdorff uniform space.

PROPOSITION 11. Let A be a dense subset of a topological space X, and let f be a mapping of A into a complete Hausdorff uniform space X'. Then f can be extended by continuity to X if and only if, for each x ∈ X, the image under f of the trace on A of the neighbourhood filter of x in X is a Cauchy filter base in X'.

This follows from the theorem of extension by continuity (loc. cit.) because \(X'\) is regular (§ 1, no. 2, Proposition 3) and because on \(X'\) convergent filters are the same as Cauchy filters.

When X is also a uniform space, there is the following theorem:

THEOREM 2. Let \(f\) be a function defined on a dense subspace \(A\) of a uniform space \(X\) , taking its values in a complete Hausdorff uniform space \(X'\) , and suppose that \(f\) is uniformly continuous on \(A\) . Then \(f\) can be extended to the whole of \(X\) by continuity, and the extended function \(\overline{f}\) is uniformly continuous.

The existence of \(\overline{f}\) is an immediate consequence of Propositions 3 and 11 of no. 1. Let us show that \(\overline{f}\) is uniformly continuous. Let \(V'\) be a closed symmetric entourage of \(X'\) , and let \(V\) be an entourage of \(X\) such that, when \(x\) and \(y\) are in \(A\) and are \(V\) -close, then \(f(x)\) and \(f(y)\)

are V'-close. We may assume that V is the closure in X × X of an entourage W of A (§ 2, no. 4, Proposition 6). We have {[} \(\overline{f}(x)\) , \(\overline{f}(y)\) {]} ∈ V' when \((x,y) \in \mathbb{W}\) ; since \(\overline{f} \times \overline{f}\) is continuous in X × X (Chapter I, § 4, no. 1, Proposition 1) we have also {[} \(\overline{f}(x)\) , \(\overline{f}(y)\) {]} ∈ V' when \((x,y) \in \mathbb{V} = \overline{\mathbb{W}}\) , since V' is closed (Chapter I, § 2, no. 1, Theorem 1).

Q.E.D.

COROLLARY. Let \(X_{1}\) , \(X_{2}\) be two complete Hausdorff uniform spaces, and let \(Y_{1}\) , \(Y_{2}\) be dense subspaces of \(X_{1}\) , \(X_{2}\) respectively. Then every isomorphism f of \(Y_{1}\) onto \(Y_{2}\) extends to an isomorphism of \(X_{1}\) onto \(X_{2}\) .

\(f\) is uniformly continuous in \(\mathbf{Y}_1\) , hence (Theorem 2) extends to a uniformly continuous mapping \(\overline{f}\colon \mathbf{X}_1\to \mathbf{X}_2\) . Likewise the inverse \(g\) of \(f\) extends to a uniformly continuous mapping \(\overline{f}^1\colon \mathbf{X}_2\to \mathbf{X}_1\) . The function \(\overline{g}\circ \overline{f}\) is therefore a continuous mapping of \(\mathbf{X}_1\) into itself whose restriction to \(\mathbf{Y}_1\) is the identity mapping; by the principle of extension of identities (Chapter I, § 8, no. 1, Corollary 1 to Proposition 2) \(\overline{g}\circ \overline{f}\) is therefore the identity mapping of \(\mathbf{X}_1\) ; similarly \(\overline{f}\circ \overline{g}\) is the identity map of \(\mathbf{X}_2\) . Consequently (Set Theory, R, § 2, no. 12) \(\overline{f}^1\) and \(\overline{g}\) are bijections and are inverses of each other; they are also uniformly continuous and are therefore isomorphisms (§ 2, no. 1, Proposition 2).

\pandocbounded{\includegraphics[keepaspectratio]{images/part002_406c0853.jpg}}

It should be remarked that if \(f\) is a bijective uniformly continuous mapping of \(\mathbf{Y}_1\) onto \(\mathbf{Y}_2\) , its extension by continuity need be neither injective nor surjective (Exercise 3).
% semantic-fragments: legacy-input-seam
\relax{}
